% Part: lambda-calculus
% Chapter: representability
% Section: currying

\documentclass[../../../include/open-logic-section]{subfiles}

\begin{document}

\olfileid{lam}{rep}{cur} 
\olsection{التكريع}

مجرّد~$\lambd$، أي~$\lambd[{\OLId{latin-ordinary}{x}}][{\OLId{latin-looped}{M}}]$، يمثل دالة لا تأخذ إلا وسيطًا
واحدًا. وقد يضيق هذا حين نطلب دالة ذات وسائط عدة.
ومن سبيل رفع القيد أن نوسع حساب~$\lambd$ فنجيز فيه
الأزواج والثلاثيات ونحوها، فتأخذ الدالة الثلاثية الموضع مثلًا،
الممثلة بـ$\lambd[{\OLId{latin-ordinary}{x}}][{\OLId{latin-looped}{M}}]$، ثلاثيةً وسيطًا لها.
غير أن \emph{التكريع} طريق أيسر إلى هذا الغرض.

نوضح ذلك بمثال، فنفترض مؤقتًا وجود عملية~$+$ في
حساب~$\lambd$. الجمع دالة ذات~${\OLMathNumeral{2}}$ موضع، تأخذ وسيطين؛
لكن مجرّد~$\lambd$ لا ينشئ إلا دالة أحادية الموضع،
ولا يجيز التركيب تعبيرًا كـ$\lambd[({\OLId{latin-ordinary}{x}},{\OLId{latin-ordinary}{y}})][({\OLId{latin-ordinary}{x}}+{\OLId{latin-ordinary}{y}})]$.
فلننظر بدلًا من ذلك في الدالة الأحادية الموضع~${\OLId{latin-ordinary}{f}}_{\OLId{latin-ordinary}{x}}({\OLId{latin-ordinary}{y}})$
التي يمثلها~$\lambd[{\OLId{latin-ordinary}{y}}][({\OLId{latin-ordinary}{x}}+{\OLId{latin-ordinary}{y}})]$، فتزيد~${\OLId{latin-ordinary}{x}}$ على وسيطها
الوحيد~${\OLId{latin-ordinary}{y}}$. وهذه في الحقيقة عائلة دوال، لا دالة واحدة:
«أضف~${\OLId{latin-ordinary}{x}}$» دالة لكل عدد~${\OLId{latin-ordinary}{x}}$.
ثم نعرّف دالة أحادية الموضع~${\OLId{latin-ordinary}{g}}$ بالحد~$\lambd[{\OLId{latin-ordinary}{x}}][{\OLId{latin-ordinary}{f}}_{\OLId{latin-ordinary}{x}}]$.
إذا أعطيناها الوسيط~${\OLId{latin-ordinary}{x}}$، كانت قيمة~${\OLId{latin-ordinary}{g}}({\OLId{latin-ordinary}{x}})$ هي الدالة~${\OLId{latin-ordinary}{f}}_{\OLId{latin-ordinary}{x}}$؛
فقيمها إذن دوال. فإذا طبقنا~${\OLId{latin-ordinary}{g}}$ على~${\OLId{latin-ordinary}{x}}$،
ثم طبقنا الحاصل على~${\OLId{latin-ordinary}{y}}$، حصل~$({\OLId{latin-ordinary}{g}}({\OLId{latin-ordinary}{x}})){\OLId{latin-ordinary}{y}} = {\OLId{latin-ordinary}{f}}_{\OLId{latin-ordinary}{x}}({\OLId{latin-ordinary}{y}}) = {\OLId{latin-ordinary}{x}}+{\OLId{latin-ordinary}{y}}$.
فتؤدي~${\OLId{latin-ordinary}{g}}$، وهي أحادية الموضع، عمل دالة الجمع الثنائية الموضع.
وهذا التحويل من دالة ثنائية الموضع إلى دالة أحادية الموضع
قيمها دوال أحادية الموضع هو ما نعنيه بالتكريع.

ولنأخذ الآن مثالًا من التركيب الصحيح لحساب~$\lambd$ نفسه.
نريد تمثيل الدالة~${\OLId{latin-ordinary}{f}}({\OLId{latin-ordinary}{x}},{\OLId{latin-ordinary}{y}}) = {\OLId{latin-ordinary}{x}}$، التي تأخذ وسيطين وتعيد
أولهما. الحد~$\lambd[{\OLId{latin-ordinary}{x}}][\lambd[{\OLId{latin-ordinary}{y}}][{\OLId{latin-ordinary}{x}}]]$ يفي بذلك:
فهو يقبل وسيطًا~${\OLId{latin-ordinary}{x}}$ ويعيد الدالة~$\lambd[{\OLId{latin-ordinary}{y}}][{\OLId{latin-ordinary}{x}}]$.
وهذه الدالة~$\lambd[{\OLId{latin-ordinary}{y}}][{\OLId{latin-ordinary}{x}}]$ تقبل وسيطًا ثانيًا~${\OLId{latin-ordinary}{y}}$،
ثم لا تعبأ به وتعيد~${\OLId{latin-ordinary}{x}}$ في كل حال.
فإذا طبقنا~$\lambd[{\OLId{latin-ordinary}{x}}][\lambd[{\OLId{latin-ordinary}{y}}][{\OLId{latin-ordinary}{x}}]]$ على وسيطين جرى الاختزال:
\begin{align*}
  (\lambd[{\OLId{latin-ordinary}{x}}][\lambd[{\OLId{latin-ordinary}{y}}][{\OLId{latin-ordinary}{x}}]])MN 
  \bredone & (\lambd[{\OLId{latin-ordinary}{y}}][{\OLId{latin-looped}{M}}]){\OLId{latin-looped}{N}} \\
  \bredone & {\OLId{latin-looped}{M}}
\end{align*}

وعلى العموم، إذا أردنا دالة ذات معلمات~${\OLId{latin-ordinary}{x}}_{\OLMathNumeral{1}}$، \dots،~${\OLId{latin-ordinary}{x}}_{\OLId{latin-ordinary}{n}}$
يحددها الحد~${\OLId{latin-looped}{N}}$، كتبنا
$\lambd[{\OLId{latin-ordinary}{x}}_{\OLMathNumeral{1}}][\lambd[{\OLId{latin-ordinary}{x}}_{\OLMathNumeral{2}}][\ldots\lambd[{\OLId{latin-ordinary}{x}}_{\OLId{latin-ordinary}{n}}][{\OLId{latin-looped}{N}}]]]$.
وبتطبيقها على~${\OLId{latin-ordinary}{n}}$ من الوسائط يكون:
\begin{multline*}
  (\lambd[{\OLId{latin-ordinary}{x}}_{\OLMathNumeral{1}}][\lambd[{\OLId{latin-ordinary}{x}}_{\OLMathNumeral{2}}][\ldots\lambd[{\OLId{latin-ordinary}{x}}_{\OLId{latin-ordinary}{n}}][{\OLId{latin-looped}{N}}]]]) {\OLId{latin-looped}{M}}_{\OLMathNumeral{1}} \dots {\OLId{latin-looped}{M}}_{\OLId{latin-ordinary}{n}} \bredone\\
  \begin{aligned}
  \bredone {} & (\Subst{(\lambd[{\OLId{latin-ordinary}{x}}_{\OLMathNumeral{2}}][\ldots\lambd[{\OLId{latin-ordinary}{x}}_{\OLId{latin-ordinary}{n}}][{\OLId{latin-looped}{N}}]])}{{\OLId{latin-looped}{M}}_{\OLMathNumeral{1}}}{{\OLId{latin-ordinary}{x}}_{\OLMathNumeral{1}}}) {\OLId{latin-looped}{M}}_{\OLMathNumeral{2}}
  \dots {\OLId{latin-looped}{M}}_{\OLId{latin-ordinary}{n}}\\
   \eqs {} & (\lambd[{\OLId{latin-ordinary}{x}}_{\OLMathNumeral{2}}][\ldots\lambd[{\OLId{latin-ordinary}{x}}_{\OLId{latin-ordinary}{n}}][\Subst{{\OLId{latin-looped}{N}}}{{\OLId{latin-looped}{M}}_{\OLMathNumeral{1}}}{{\OLId{latin-ordinary}{x}}_{\OLMathNumeral{1}}}]]) {\OLId{latin-looped}{M}}_{\OLMathNumeral{2}}
            \dots {\OLId{latin-looped}{M}}_{\OLId{latin-ordinary}{n}} \\
  \vdots & \\
  \bredone {} &\Subst{\Subst{{\OLId{latin-looped}{N}}}{{\OLId{latin-looped}{M}}_{\OLMathNumeral{1}}}{{\OLId{latin-ordinary}{x}}_{\OLMathNumeral{1}}}\ldots}{{\OLId{latin-looped}{M}}_{\OLId{latin-ordinary}{n}}}{{\OLId{latin-ordinary}{x}}_{\OLId{latin-ordinary}{n}}}
  \end{aligned}
\end{multline*}
فمؤدى السطر الأخير إحلال~${\OLId{latin-looped}{M}}_{\OLId{latin-ordinary}{i}}$ محل~${\OLId{latin-ordinary}{x}}_{\OLId{latin-ordinary}{i}}$ في متن تعريف
الدالة؛ وهذا هو المطلوب من تطبيقها على وسائط متعددة.
\end{document}
